% =============================================================================
%  config/macros.tex: comandi semantici
%  Nei file di contenuto usare SEMPRE questi comandi: se si vuole cambiare la
%  notazione basta modificarla qui.
% =============================================================================

% --- Insiemi numerici -----------------------------------------------------------
\newcommand{\R}{\mathbb{R}}
\newcommand{\Z}{\mathbb{Z}}
\newcommand{\N}{\mathbb{N}}

% --- Algebra lineare -----------------------------------------------------------------
\newcommand{\transp}{^{\mathsf{T}}}                  % trasposto:  c\transp x
\newcommand{\inv}{^{-1}}                             % inversa:    A_B\inv
\DeclarePairedDelimiter{\norm}{\lVert}{\rVert}       % \norm{x}
\DeclarePairedDelimiter{\abs}{\lvert}{\rvert}        % \abs{x}
\DeclarePairedDelimiter{\floor}{\lfloor}{\rfloor}    % \floor{x}  (tagli di Gomory)
\DeclarePairedDelimiter{\ceil}{\lceil}{\rceil}       % \ceil{x}

% --- Operatori -------------------------------------------------------------------------
\DeclareMathOperator*{\argmin}{arg\,min}
\DeclareMathOperator*{\argmax}{arg\,max}
\DeclareMathOperator{\conv}{conv}                    % inviluppo convesso
\DeclareMathOperator{\cone}{cono}                    % cono generato
\DeclareMathOperator{\rank}{rango}

% --- Notazione PL / PLI ------------------------------------------------------------------
% Problema primale e duale in forma standard del libro:
%   (P)  max { c^T x : A x <= b }        (D)  min { y^T b : y^T A = c^T, y >= 0 }
\newcommand{\primale}{(\mathcal{P})}
\newcommand{\duale}{(\mathcal{D})}
\newcommand{\vP}{v(\mathcal{P})}                     % valore ottimo del primale
\newcommand{\vD}{v(\mathcal{D})}                     % valore ottimo del duale
\newcommand{\regP}{\Omega}                           % regione ammissibile
\newcommand{\vs}{v_S}                                % valutazione superiore (B&B)
\newcommand{\vi}{v_I}                                % valutazione inferiore (B&B)

% Problema di ottimizzazione su più righe:
%   \problema{\max}{c\transp x}{A x \le b \\ x \in \Z^n}
\newcommand{\problema}[3]{%
  \begin{cases}
    #1\ #2 \\
    \begin{aligned}[t] #3 \end{aligned}
  \end{cases}}

% --- Testo ----------------------------------------------------------------------------------
\newcommand{\termine}[1]{\textbf{#1}}                % termine definito per la prima volta
\newcommand{\TODO}[1]{\textcolor{red}{\textbf{[TODO: #1]}}}
